\section{题05：生态廊道（P16787）}

\subsection{问题建模：路径图上局部重排的边权最大化}

题目给出了$n$个排成一列的监测位点，每个位点$i$拥有一个基因多样性编码$a_i$（$0\le a_i\le 10^9$）。定义交互强度为所有相邻位点编码按位异或之和：$S=\sum_{i=1}^{n-1}(a_i\oplus a_{i+1})$。我们可以将这条序列建模为一条含有$n$个顶点的路径图，第$i$条边连接顶点$i$和$i+1$，其权值即为$a_i\oplus a_{i+1}$，总权值$S$便是交互强度。

工程师的优化操作是：在全序列上选择两个互不重叠、长度均为$3$的连续区域（只能选择一次），对每个区域内部的三个位点进行任意重新排列。排列只改变数值在区域内的出现顺序，不改变数值集合本身。记两个区域的起点下标分别为$i$和$j$（$1\le i,j\le n-2$），约束$|i-j|\ge 3$保证了两个长度为$3$的区域不会交叉或重叠。我们的目标是在所有满足条件的$(i,j)$以及所有可能的排列中，最大化最终的交互强度。

为了与代码保持一致，下文统一采用0-indexed下标（数组下标范围为$[0,n-1]$）。此时区间起点$i$（$0\le i\le n-3$）所覆盖的三个位置为$\{i,i+1,i+2\}$。可选的合法起点总数为$m=n-2$。原始交互强度$S_{\text{base}}=\sum_{k=0}^{n-2}(a_k\oplus a_{k+1})$在输入后即可$\mathrm{O}(n)$求出，所有优化方案的总强度均可写为$S_{\text{base}}+\Delta$的形式，其中$\Delta$为排列操作带来的总增量。于是问题等价于求$\Delta$的最大值。

\subsection{朴素暴力：从三次方到局部增量的第一步收缩}

最直接的暴力解法分三层：第一层枚举区间起点对$(i,j)$，共有$\mathrm{O}(n^2)$种有效组合；第二层，每个区间有$3!=6$种排列，组合起来共$36$种方案；第三层，对每种方案重新计算全部$n-1$条边的异或和，耗时$\mathrm{O}(n)$。总复杂度$\mathrm{O}(36n^3)$，当$n=2\times10^5$时$n^3\approx8\times10^{15}$，完全不可行。

稍加观察即可发现优化空间：一个长度为$3$的区间被重新排列后，受影响的边总共不超过四条——区间内部有两条边$(i,i+1)$和$(i+1,i+2)$，区间左侧最多有一条跨越边界的边$(i-1,i)$（仅当$i>0$时存在），区间右侧最多也有一条跨越边界的边$(i+2,i+3)$（仅当$i+2<n-1$即$i<n-3$时存在）。序列中其余所有边的贡献均不发生变化。因此对每组$(i,j,p,q)$，我们可以仅重算这至多$8$条受影响边的贡献变化，在$\mathrm{O}(1)$时间内完成增量的计算，整体复杂度降为$\mathrm{O}(36n^2)$。然而$n=2\times10^5$时$n^2\approx4\times10^{10}$，依然远超时限。问题的关键在于：我们不应该先确定$(i,j)$再计算增量，而应当反过来——先独立计算出每个区间各自的最优增量，再利用结构性质高效地合并它们。

\subsection{核心洞察：区间增量可独立预计算，间距决定合并方式}

考察单个区间：设0-indexed下标$i$为起点，原值$(a_i,a_{i+1},a_{i+2})=(x,y,z)$。该区间的某种排列$p$将三个位置上的值分别变为$b_0,b_1,b_2$（三者是$(x,y,z)$的一个置换）。该排列对交互强度的增量$\Delta_i(p)$由三部分组成：
\[
\begin{aligned}
\Delta_i^{\text{left}}(p) &=
\begin{cases}
(a_{i-1}\oplus b_0)-(a_{i-1}\oplus x), & i>0,\\
0, & i=0;
\end{cases}\\[4pt]
\Delta_i^{\text{in}}(p) &= (b_0\oplus b_1)+(b_1\oplus b_2)-(x\oplus y)-(y\oplus z);\\[4pt]
\Delta_i^{\text{right}}(p) &=
\begin{cases}
(b_2\oplus a_{i+3})-(z\oplus a_{i+3}), & i<n-3,\\
0, & i=n-3.
\end{cases}
\end{aligned}
\]
三者之和即为区间$i$在排列$p$下的总增量$\Delta_i(p)$。对每个起点$i$，枚举全部$6$种排列即可在$\mathrm{O}(1)$时间内得到该处的最优独立增量$\Delta_i^{\max}=\max_p\Delta_i(p)$。

这些“独立增量”的隐含假设是：区间两端的邻居位点均保持原始值不变。当两个所选区间距离足够远（中间至少隔了一个未被调整的位置），这一假设对两个区间同时成立——它们各自的受影响边集合互不相交，总增量可以直接相加：$\Delta=\Delta_i(p)+\Delta_j(q)$。然而，若两个区间恰好紧邻（即$j=i+3$），则区间$i$的右边界和区间$j$的左边界恰好共享同一条边$(i+2,i+3)$，两套独立增量都试图修改这条边、却各自假设对端为原值，从而导致叠加出错。这一观察自然将问题划分为两种情形：独立情形（$|i-j|\ge 4$）和共享边界情形（$|i-j|=3$）。

\subsection{独立情形：前缀后缀最大值将配对从O(n²)压缩至O(n)}

当$|i-j|\ge 4$时，不妨设$j\ge i+4$。由于区间$i$覆盖$[i,i+2]$，区间$j$覆盖$[j,j+2]$，位置$i+3$位于两者之间且未被任何区间修改——它恰好是区间$i$右边界所使用的“原始邻居”，也是区间$j$左边界所依赖的那个位置的左侧元素。两个区间的受影响边集合完全不相交，故总最优增量为$\Delta_i^{\max}+\Delta_j^{\max}$。问题转化为：在所有满足$|i-j|\ge 4$的起点对中，求$\Delta_i^{\max}+\Delta_j^{\max}$的最大值。

这一子问题通过前缀最大值和后缀最大值在$\mathrm{O}(m)$时间内解决。定义前缀数组$\mathit{pref}[k]=\max_{0\le t\le k}\Delta_t^{\max}$和后缀数组$\mathit{suff}[k]=\max_{k\le t<m}\Delta_t^{\max}$，均为$\mathrm{O}(m)$建表。对每个起点$i$，所有满足$j\le i-4$的$j$中最优者即为$\mathit{pref}[i-4]$（当$i\ge4$时有效），所有满足$j\ge i+4$的$j$中最优者即为$\mathit{suff}[i+4]$（当$i+4<m$时有效）。遍历全部$i$，检查$\Delta_i^{\max}+\mathit{pref}[i-4]$和$\Delta_i^{\max}+\mathit{suff}[i+4]$，取全局最大值即覆盖了所有独立情形的最优配对。

以样例验证：$n=7$，$a=[1,2,3,4,5,6,7]$（1-indexed），则$m=5$个0-indexed起点的原始编码分别为：$i=0$覆盖$[1,2,3]$，$i=1$覆盖$[2,3,4]$，$i=2$覆盖$[3,4,5]$，$i=3$覆盖$[4,5,6]$，$i=4$覆盖$[5,6,7]$。对于$i=1$，排列$[2,4,3]$给出内部增量$(2\oplus4)+(4\oplus3)-(2\oplus3)-(3\oplus4)=6+7-1-7=5$，右边界增量$(3\oplus5)-(4\oplus5)=6-1=5$，故$\Delta_1^{\max}=10$；$i=4$的最优排列即原序，$\Delta_4^{\max}=0$。注意$|1-4|=3$，恰为共享边界情形，独立情形的前缀/后缀扫描不会覆盖它。

\subsection{共享边界：精确修正重叠边的重复调整}

当$|i-j|=3$时，由对称性只需考虑$j=i+3$（另一方向在枚举中自动覆盖）。区间$i$覆盖$[i,i+1,i+2]$，区间$j$覆盖$[i+3,i+4,i+5]$，两者仅通过边$(i+2,i+3)$相连。记区间$i$在排列$p$下的右端值为$R_i(p)$（即原来的$b_2$），区间$j$在排列$q$下的左端值为$L_j(q)$（即$c_0$），原始共享边的贡献$\mathit{orig}=a_{i+2}\oplus a_{i+3}$。直接相加$\Delta_i(p)+\Delta_j(q)$会将同一原始贡献扣减两次、又将各自基于“对端不变”假设的新贡献各加一次，形成偏差。正确的合并增量需要先卸掉两区间各自对共享边的独立修正，再补上两个新区间端点直接异或的真实贡献：
\[
\Delta_{\text{combined}}(p,q)=\Delta_i(p)+\Delta_j(q)-\Delta_i^{\text{right}}(p)-\Delta_j^{\text{left}}(q)+(R_i(p)\oplus L_j(q))-\mathit{orig}.
\]

回到样例验证此公式。取$i=1$、$j=4$，$R_1(p)=3$（排列$[2,4,3]$的右端值），$L_4(q)=5$（原始排列的左端值），$\mathit{orig}=a_3\oplus a_4=4\oplus5=1$。已知$\Delta_1(p)=10$、$\Delta_1^{\text{right}}(p)=5$，$\Delta_4(q)=0$、$\Delta_4^{\text{left}}(q)=0$。代入得$10+0-5-0+(3\oplus5)-1=5+6-1=10$。加上$S_{\text{base}}=16$后总强度为$26$，与样例答案一致。

对于每一对满足$j=i+3$的起点，枚举两个区间各自的$6$种排列（共$36$种组合），用上述公式计算合并增量并维护全局最大值。值得注意的是，这里不能简单地取$\Delta_i^{\max}+\Delta_j^{\max}$再作边界修正——因为使$\Delta_i$最大的那个排列$p$，其右端值$R_i(p)$不一定与使$\Delta_j$最大的那个排列$q$的左端值$L_j(q)$搭配时得到最大的$(R_i\oplus L_j)$，所以必须遍历全部$36$种组合。好在这一枚举每对只需常数时间，总开销$\mathrm{O}(36n)$完全可接受。

\subsection{整体复杂度与实现要点}

算法分为三个阶段，均为线性。预处理阶段：对$m=n-2$个起点各枚举$6$种排列，每个排列在$\mathrm{O}(1)$时间内完成三条增量的计算，总耗时$\mathrm{O}(6m)=\mathrm{O}(n)$。独立情形阶段：构建$\mathit{pref}$和$\mathit{suff}$数组各需一次$\mathrm{O}(m)$扫描，配对扫描又为$\mathrm{O}(m)$。共享边界阶段：枚举$\mathrm{O}(n)$对紧邻区间，每对枚举$36$种排列组合，总耗时$\mathrm{O}(36n)$。因此总体时间复杂度$\mathrm{O}(n)$，空间复杂度$\mathrm{O}(n)$（主要为存储每起点的$6$组增量分量和端点值）。对于$n=2\times10^5$的数据量，约$7.2\times10^6$次共享边界枚举运算，在现代CPU上毫秒级可完成。

实现时需注意几个细节。异或值的累加可能达到$2\times10^5\cdot 2^{30}\approx 2\times10^{14}$，超出的\mil{int}的$2^{31}-1\approx2.1\times10^9$，必须使用64位有符号整数（C++中的\mil{long long}）。预处理时需为每个起点的每种排列同时保存总增量\mil{total_delta}、左边界增量\mil{delta_left}、右边界增量\mil{delta_right}以及左右端点实际值\mil{left_val}和\mil{right_val}，供共享边界阶段取用。边界位置的增量为零需显式判定：起点$i=0$时左边界分量记为$0$，$i=m-1$（即区间右端已触碰序列末尾）时右边界分量记为$0$。共享边界情形中，将$a_{i+2}$与$a_{i+3}$的异或值作为$\mathit{orig}$直接代入公式即可，无需引入额外的索引映射。

\subsection{代码实现}
\inputminted{c++}{../src/p05_ecology_corridor.cpp}
